\section{Conclusion}
\label{sec:conclusion}

We have presented DeepChrInteract-v2, a modern reimplementation of the original
DeepChrInteract toolkit that brings the codebase to current PyTorch standards while
substantially expanding its scientific scope.
The key engineering contribution is the elimination of the PNG-based preprocessing
pipeline in favour of in-memory on-the-fly one-hot encoding, reducing preprocessing
time from hours to seconds and removing gigabytes of intermediate disk storage.

On the scientific side, four new encoder families---BiLSTM, Transformer,
CNN+BiLSTM hybrid, and CNN+Transformer hybrid---extend the coverage of the original
benchmark to include the sequence modelling paradigms dominant in contemporary
computational genomics.
The integration of DNA foundation models (DNABERT-2, Nucleotide Transformer, HyenaDNA)
as plug-in feature extractors positions the framework to leverage the rapid progress
in genomic pre-training, while the configurable multi-path fusion module enables
fine-grained study of enhancer--promoter interaction modelling strategies.

\paragraph{Limitations.}
(i) Due to the absence of GPU-accelerated servers during development, quantitative
results in \Cref{sec:experiments} are placeholders pending experimental completion.
(ii) LLM fine-tuning experiments are computationally expensive and may require
$>$24~GB GPU memory for the largest Nucleotide Transformer variant.
(iii) The dataset is limited to GM12878; broader cell-type benchmarks
(e.g., K562, IMR90) are planned for future work.

\paragraph{Future Directions.}
\begin{itemize}[noitemsep]
  \item \textbf{Cross-attention fusion.}
        Rather than post-hoc vector operations, a cross-attention layer between
        enhancer and promoter token sequences may capture interaction-specific
        motif co-occurrences more faithfully.
  \item \textbf{Contrastive pre-training.}
        Self-supervised contrastive objectives on matched E--P pairs from
        Hi-C data could provide richer representations without relying solely
        on the limited number of labelled interactions.
  \item \textbf{Multi-species transfer.}
        Given the cross-species pre-training of DNABERT-2 and Nucleotide Transformer,
        it would be valuable to assess whether models trained on human data generalise
        to orthologous loci in mouse or zebrafish.
\end{itemize}

\paragraph{Reproducibility.}
All code, configuration files, and preprocessing scripts are available at
\url{https://github.com/lichen-lab/DeepChrInteract} (original) and in the
\texttt{DeepChrInteract-v2/} directory of the same repository.
All random seeds, hyperparameters, and evaluation protocols are specified in
\texttt{src/config.py} to ensure full reproducibility.
